\chapter{Data Logging and Analysis Workflow}
\label{app:data_logging}

This appendix explains the complete pipeline from recording measurements
on the robot to producing publication-quality plots on any computer.
The same pattern is used in all four experiments.

%% ─────────────────────────────────────────────────────────────────
\section{Overview}
%% ─────────────────────────────────────────────────────────────────

\begin{figure}[htb]
\centering
\begin{tikzpicture}[
  box/.style={draw, rounded corners=3pt, minimum width=4.5cm,
              minimum height=0.85cm, align=center, font=\small},
  arr/.style={-{Latex}, thick}
]
  \node[box, fill=blue!10]   (robot) at (0,0)    {Franka Panda\\(Robot PC, RT-Linux)};
  \node[box, fill=orange!15] (cpp)   at (0,-1.8) {C++ Controller\\(1 ms callback)};
  \node[box, fill=green!10]  (csv)   at (0,-3.6) {CSV File\\(\texttt{exp1\_data.csv})};
  \node[box, fill=red!10]    (pc)    at (0,-5.4) {Analysis Machine\\(any \abbr{OS}, no robot)};
  \node[box, fill=purple!10] (plots) at (0,-7.2) {PDF Plots\\(\texttt{plot\_exp1.py})};

  \draw[arr] (robot) -- (cpp)   node[midway,right,font=\tiny]{state at 1 kHz};
  \draw[arr] (cpp)   -- (csv)   node[midway,right,font=\tiny]{append row per ms};
  \draw[arr] (csv)   -- node[midway,right,font=\small]{\texttt{scp}} (pc);
  \draw[arr] (pc)    -- (plots) node[midway,right,font=\tiny]{\texttt{pandas.read\_csv()}};
\end{tikzpicture}
\caption{Data logging and analysis pipeline.
         Measurements are written to CSV during the experiment, copied
         to an analysis machine via \texttt{scp}, and plotted with Python.
         The robot PC and the analysis machine can be completely separate.}
\label{fig:data_pipeline}
\end{figure}

%% ─────────────────────────────────────────────────────────────────
\section{Step 1 --- Writing CSV in C++ (Robot PC)}
%% ─────────────────────────────────────────────────────────────────

The C++ controller opens a CSV file before the control loop starts,
writes one row per millisecond inside the callback, and closes the file
after the loop ends.
No additional libraries are required --- standard C++ \texttt{<fstream>}
is used.

\begin{lstlisting}[style=cppstyle,
  caption={C++ CSV logging pattern used in all experiments},
  label={lst:csv_write}]
#include <fstream>   // standard library -- no extra dependency

// -- BEFORE the control loop: open file and write header --------
std::ofstream csv("exp1_data.csv");   // file created in current dir
csv << "time,e_px,e_py,e_pz,"        // header row
       "fx,fy,fz,"
       "tau1,tau2,tau3,tau4,tau5,tau6,tau7\n";

// -- INSIDE the 1 ms callback: append one row -------------------
robot.control([&](const franka::RobotState& state,
    franka::Duration dt) -> franka::Torques
{
    time += dt.toSec();

    // ... compute e_p, f, tau as usual ...

    // Write one row every millisecond
    csv << time
        << "," << e_p(0) << "," << e_p(1) << "," << e_p(2)
        << "," << f(0)   << "," << f(1)   << "," << f(2);
    for (int i = 0; i < 7; ++i) csv << "," << tau(i);
    csv << "\n";          // flush happens automatically at close

    // ... return torques as usual ...
});

// -- AFTER the control loop: close file -------------------------
csv.close();
std::cout << "Data saved to exp1_data.csv\n";
\end{lstlisting}

\subsection{CSV File Format}

Each row contains one measurement at one time step (1\,ms interval).
The first row is a header with column names.
A typical file after a 5\,s experiment contains $5{,}000$ rows and is
approximately 300\,KB --- small enough to transfer instantly.

\begin{center}
\renewcommand{\arraystretch}{1.3}
\begin{tabular}{@{}llll@{}}
\toprule
Column & Type & Unit & Description \\
\midrule
\texttt{time}     & float & s   & Time since controller start \\
\texttt{e\_px/y/z} & float & m   & Translational error $e_p$ \\
\texttt{p\_x/y/z} & float & m   & End-effector position $p_{EE}$ \\
\texttt{phi\_err}  & float & rad & Rotation error magnitude $\phi$ \\
\texttt{fx/y/z}   & float & N   & Cartesian restoring force \\
\texttt{mx/y/z}   & float & Nm  & Cartesian restoring moment \\
\texttt{tau1--7}  & float & Nm  & Joint torques \\
\texttt{tau\_g1--7}& float & Nm  & Gravity compensation torques \\
\texttt{F\_ext\_x/y/z}& float & N & libfranka external force estimate \\
\texttt{d\_wall\_mm}  & float & mm & Signed distance from wall \\
\texttt{phi\_tilt\_deg}& float & deg & Tilt angle (free \abbr{DOF}) \\
\texttt{phi\_spin\_deg}& float & deg & Spin angle (stiff \abbr{DOF}) \\
\bottomrule
\end{tabular}
\end{center}

%% ─────────────────────────────────────────────────────────────────
\section{Step 2 --- Copying Files to the Analysis Machine}
%% ─────────────────────────────────────────────────────────────────

After the experiment, copy the CSV files from the robot PC to any
machine (laptop, desktop) using \texttt{scp}:

\begin{lstlisting}[style=cppstyle, language=bash,
  caption={Copy CSV files from the robot workstation},
  label={lst:scp}]
# From the analysis machine (not the robot PC):
scp ubuntu@172.16.0.1:~/exp1_data.csv  ./data/
scp ubuntu@172.16.0.1:~/exp2_data.csv  ./data/
scp ubuntu@172.16.0.1:~/exp3_*.csv     ./data/
scp ubuntu@172.16.0.1:~/exp4_*.csv     ./data/

# Or copy all CSV files at once:
scp ubuntu@172.16.0.1:~/*.csv          ./data/
\end{lstlisting}

The analysis machine does not need libfranka, a real-time kernel,
or any connection to the robot.
It only needs Python with pandas and matplotlib.

%% ─────────────────────────────────────────────────────────────────
\section{Step 3 --- Reading and Plotting in Python}
%% ─────────────────────────────────────────────────────────────────

\subsection{Installation}

On the analysis machine (once only):

\begin{lstlisting}[style=cppstyle, language=bash,
  caption={Install Python analysis dependencies},
  label={lst:pip}]
pip install pandas matplotlib numpy
\end{lstlisting}

\subsection{Reading a CSV file}

\begin{lstlisting}[style=cppstyle, language=Python,
  caption={Loading and inspecting a CSV file with pandas},
  label={lst:csv_read}]
import pandas as pd
import numpy as np
import matplotlib.pyplot as plt

# Load the CSV -- pandas infers column names from the header row
df = pd.read_csv("data/exp1_data.csv")

# Quick inspection
print(df.shape)          # (5000, 21) -- 5000 rows, 21 columns
print(df.columns.tolist())
print(df.head(3))        # first 3 rows

# Access individual columns by name
t    = df["time"].values          # time vector [s]
e_px = df["e_px"].values * 1000   # convert m -> mm
f_x  = df["fx"].values            # Cartesian force [N]
tau1 = df["tau1"].values          # joint 1 torque [Nm]
\end{lstlisting}

\subsection{Plotting from CSV data}

\begin{lstlisting}[style=cppstyle, language=Python,
  caption={Plotting position error and joint torques from measured data},
  label={lst:csv_plot}]
fig, (ax1, ax2) = plt.subplots(2, 1, figsize=(10, 7), sharex=True)

# Plot position error (x-axis)
ax1.plot(t, df["e_px"]*1000, color="#1f77b4", linewidth=2,
         label="$e_{p,x}(t)$  (measured)")
ax1.axhline(0, color="k", linewidth=0.7, linestyle="--")
ax1.set_ylabel("$e_{p,x}$ [mm]")
ax1.set_title("Experiment 1 -- Measured Step Response")
ax1.legend()

# Plot joint torques
colors = plt.cm.tab10(range(7))
for i in range(1, 8):
    ax2.plot(t, df[f"tau{i}"], color=colors[i-1],
             linewidth=1.5, label=f"$\\tau_{i}$")
ax2.set_ylabel("Joint torque [Nm]")
ax2.set_xlabel("Time $t$ [s]")
ax2.legend(ncol=4)

plt.tight_layout()
plt.savefig("exp1_measured.pdf", bbox_inches="tight", dpi=200)
plt.show()
\end{lstlisting}

\subsection{Comparing simulation with measurement}

\begin{lstlisting}[style=cppstyle, language=Python,
  caption={Overlay simulated prediction on measured data},
  label={lst:compare}]
# Simulated prediction (analytical)
omega_n = np.sqrt(200.0)   # K=200 N/m, M=1 kg
e_sim   = 0.050 * (1 + omega_n*t) * np.exp(-omega_n*t) * 1000  # [mm]

# Measured data from CSV
e_meas  = df["e_px"].values * 1000  # [mm]

fig, ax = plt.subplots(figsize=(9, 4))
ax.plot(t, e_sim,  "--", color="gray",    linewidth=2, label="Simulation")
ax.plot(t, e_meas, "-",  color="#1f77b4", linewidth=2, label="Measured")
ax.set_xlabel("Time $t$ [s]")
ax.set_ylabel("$e_{p,x}$ [mm]")
ax.set_title("Simulation vs. Measured Step Response  ($K=200$ N/m)")
ax.legend()
plt.tight_layout()
plt.savefig("exp1_sim_vs_meas.pdf", bbox_inches="tight")
\end{lstlisting}

%% ─────────────────────────────────────────────────────────────────
\section{Complete Plot Scripts}
%% ─────────────────────────────────────────────────────────────────

The complete plotting scripts for all four experiments are provided
on the digital medium:

\begin{center}
\renewcommand{\arraystretch}{1.3}
\begin{tabular}{@{}p{3.8cm}p{3.2cm}p{5cm}@{}}
\toprule
Script & Reads & Produces \\
\midrule
\texttt{plot\_exp1.py} & \texttt{exp1\_data.csv} &
  Error, force, torques, $\phi$ \\
\texttt{plot\_exp2.py} & \texttt{exp2\_data.csv} &
  Wall overview, forces and moments, torques \\
\texttt{plot\_exp3\_exp4.py} & \texttt{exp3/4\_*.csv} &
  Stiffness \& torque plots \\
\bottomrule
\end{tabular}
\end{center}

Execute any script as:
\begin{lstlisting}[style=cppstyle, language=bash,
  caption={Execute a plot script on measured data},
  label={lst:run_plot}]
# Default: reads CSV from current directory
python3 plot_exp1.py

# Or specify a path:
python3 plot_exp1.py data/exp1_data.csv
\end{lstlisting}

%% =============================================================
%%  APPENDIX C (was B) – EXPERIMENT 1: CARTESIAN POSITION AND ORIENTATION TRACKING
%% =============================================================
